\begin{abstract}
Constelações de satélites de radar de abertura sintética (SAR) oferecem muitas oportunidades de aquisição sobre uma mesma região, e escolher quais passagens adquirir é um problema combinatório: é preciso cobrir as áreas de interesse com poucas aquisições e pouca sobreposição entre faixas. Este trabalho aplica um algoritmo genético (AG) com codificação binária à seleção de 702 passagens da constelação ICEYE sobre sete áreas de interesse na Amazônia, que somam cerca de 210~mil~km$^2$. A aptidão combina, em uma soma ponderada, a fração da área coberta, a fração de passagens selecionadas e a sobreposição entre faixas. São comparados os métodos de seleção por roleta, torneio e ranking linear, diferentes taxas de cruzamento e mutação e a execução com e sem elitismo, pelas curvas de convergência e pela aptidão final. O torneio, de maior pressão de seleção, obteve a maior aptidão entre os métodos de seleção; a taxa de mutação mais alta (0,15) impediu a evolução; e o elitismo garantiu o crescimento sem oscilações da melhor aptidão, enquanto as demais variações de taxa ficaram dentro da variação entre execuções. Nenhuma configuração convergiu em 40 gerações, o que indica a necessidade de mais gerações e de várias execuções por cenário.
\end{abstract}

\begin{IEEEkeywords}
algoritmo genético, cobertura de área, elitismo, métodos de seleção, otimização combinatória, radar de abertura sintética
\end{IEEEkeywords}
